% Report section owner: Ayushi Negi.
% AI-assisted working draft: review the wording, sources and scope with the team before submission.
\section{Problem Definition}
\label{sec:Problem Definition}

\subsection{Stakeholder and motivation}
Financial analysts reviewing Australian bank results need to locate, interpret and compare information spread across lengthy reports. Finding a figure is only part of the task. Before using it in a briefing, an analyst needs to understand what it measures, which reporting period it covers and whether it can be compared with another result. This context may be distributed across results releases, management commentary and accounting notes.

Commonwealth Bank of Australia's (CBA) FY2025 Profit Announcement illustrates the problem. It reports four profit figures: cash and statutory profit, each for continuing operations and for the total including discontinued operations. For example, continuing cash profit is AUD~10,252 million, while total statutory profit is AUD~10,116 million. Both are valid, but they represent different results. An answer that omits these distinctions can be misleading even when its number appears in the report~\cite{fintraceCBA2025}.

FinTrace addresses this problem of finding financial information without losing its meaning. It supports an analyst's review by making the answer and its evidence available together. The analyst remains responsible for interpreting the information before using it in a briefing or financial review.

\subsection{Aim and project questions}
The project aims to develop and evaluate a RAG assistant that helps financial analysts find and understand reported information. Users can ask questions, compare supported results across years, look up financial terms and explore management's explanations. Answers provide relevant context, source pages and calculation workings where applicable.

The project investigates four questions:
\begin{enumerate}
\item Can FinTrace find the requested information while preserving its company, reporting period, unit, financial meaning and the parts of the business it covers?
\item Can it compare figures measured on the same basis across two financial years within a bank and show how the change was calculated?
\item Can it explain financial terms and reported changes using relevant evidence without presenting an unsupported interpretation as a reported fact?
\item Can it recognise when the available evidence is insufficient and communicate that limitation without inventing an answer or substituting another measure?
\end{enumerate}
Evaluation considers correctness, completeness, supporting citations, calculations and the handling of unanswered questions. Declining a request can be appropriate when evidence is insufficient, but declining a question the system should answer is a limitation. Response time is assessed separately.

\subsection{Scope and changes since Milestone 1}
The prototype covers six official reports from CBA and National Australia Bank (NAB) for the financial years FY2024 and FY2025. Milestone~1 proposed an assistant for corporate reports with separate Ask and Verify modes, including claim verdicts. We refined this into one question interface focused on banking. Selecting one sector and two years made it feasible to understand the terminology, check reference answers and evaluate historical comparisons~\cite{fintraceMilestone2026,fintraceSources2026}.

Initial testing focused on 13 core financial measures and selected tables. However, treating that list as the boundary for every question excluded useful definitions and management commentary. We therefore expanded text search across the full reports and added checks for further NAB financial tables. General profit and earnings per share questions also retain the available labelled alternatives. Calculations remain restricted to values checked against their sources~\cite{fintraceCoverage2026,fintraceBroad2026}.

Supported comparisons are between years within the same bank. Comparisons between banks remain outside the prototype's scope and would require additional checks of reporting periods and definitions. CBA and NAB have different financial year ends. Live prices, future profit predictions, private account information and personalised investment advice are also excluded. The deliverable is a local research prototype that supports evidence checking rather than replacing professional financial judgement.
